% Cover letter for the IJIRA submission. Build: pdflatex cover_letter
\documentclass[11pt,a4paper]{article}
\usepackage[a4paper,margin=2.5cm]{geometry}
\usepackage[T1]{fontenc}
\usepackage{lmodern}
\usepackage[hidelinks]{hyperref}
\setlength{\parindent}{0pt}
\setlength{\parskip}{6pt}
\pagestyle{empty}

\begin{document}

\textbf{Xuan Bach Nguyen}\\
School of Electrical and Electronic Engineering, Hanoi University of Science and Technology,\\
No. 1 Dai Co Viet, Hanoi, Vietnam\\
Faculty of Law, Foreign Trade University, 91 Chua Lang Street, Hanoi, Vietnam\\
E-mail: \href{mailto:nguyenxuanbach513@gmail.com}{nguyenxuanbach513@gmail.com} \quad ORCID: 0009-0003-5194-8790

\medskip
18 October 2026

\medskip
The Editor-in-Chief\\
\emph{International Journal of Intelligent Robotics and Applications}

\medskip
Dear Editor-in-Chief,

I submit the manuscript entitled \emph{``When to run depth inference? Observed-free-space perception scheduling for
low-cost quadrotors with delayed monocular depth''} for consideration as an original research article in the
\emph{International Journal of Intelligent Robotics and Applications}.

Low-cost quadrotors often obtain obstacle information from a single camera and a monocular depth network, which is
the most computationally expensive stage of the loop and has a latency of about 0.1~s. The manuscript asks when such
a network must be run, and answers with a geometric condition: an inference is triggered only when the planned path,
extended by the distance needed to stop after the next result arrives, leaves the free space observed by the most
recent frames. Under stated assumptions, skipping an inference then does not remove the ability to stop within
observed space. The scheduler has a hard rate cap and is integrated with a multi-rate architecture (vertical-cylinder
obstacle memory with a Kalman update at the capture time, a 20~Hz model predictive controller and a 100~Hz control
barrier function filter).

The consequences of the condition are tested in 36\,000 paired simulated flights. At the same rate cap, the
scheduler uses 50--72\,\% of the inferences of fixed-rate perception in the nominal and everyday scenarios with no
statistically significant difference in collision rate; its frame rate increases with obstacle density; at the same
number of inferences it yields a lower collision rate than scheduling by the remaining safe time in the
nominal and hard scenarios; and at an equal low
computational budget, failed missions in the nominal scenarios decrease from 35 to 16 of 720. The study also shows
that a wide-angle camera is a precondition for these savings. The work links perception, planning and control of an
aerial robot under limited onboard computation, which I believe is relevant to the readers of the journal.

The manuscript is original, has not been published, and is not under consideration for publication elsewhere. I am
the sole author. No funding was received for this study, and I have no competing interests to declare. The use of an
AI assistant in preparing the code, the analysis scripts and the first draft is described in Section~5 of the
manuscript, as required by the journal's guidelines; I take full responsibility for the content. The source code,
the per-run results and the scripts that produce the tables and figures are available at
\url{https://github.com/nxbnxb1/project1}.

Thank you for considering this submission.

\medskip
Sincerely,

\medskip
Xuan Bach Nguyen

\end{document}
